\chapter{Preliminarii}

\section{Arhitectura aplicației}

Aplicația dezvoltată utilizează o arhitectură modernă bazată pe microservicii, cu separarea clară între componenta backend și frontend. Această abordare permite scalabilitatea independentă a fiecărei componente și facilitează mentenanța pe termen lung.

\section{Tehnologii Backend}

\subsection{FastAPI}

FastAPI reprezintă framework-ul web ales pentru dezvoltarea API-ului aplicației, fiind unul dintre cele mai performante framework-uri Python disponibile. Principalele avantaje care au motivat această alegere includ performanța ridicată comparabilă cu Node.js și Go, generarea automată a documentației OpenAPI, validarea automată a datelor folosind type hints-urile Python și suportul nativ pentru programarea asincronă. Această combinație de caracteristici este esențială pentru o aplicație care procesează imagini și efectuează apeluri către servicii externe de inteligență artificială.

\subsection{PostgreSQL}

PostgreSQL a fost selectat ca sistem de gestiune a bazelor de date pentru fiabilitatea sa dovedită și conformitatea completă cu standardele SQL. Alegerea se bazează pe robustețea sistemului cu garanții ACID complete, suportul pentru tipuri de date complexe precum JSON pentru stocarea informațiilor nutriționale structurate, funcționalitățile native de căutare în text pentru implementarea funcțiilor de search și capacitățile avansate de indexare pentru optimizarea performanței interogărilor.

\subsection{Tesseract OCR}

Pentru recunoașterea optică a caracterelor, Tesseract OCR oferă una dintre cele mai performante soluții open-source disponibile. Dezvoltat inițial de Hewlett-Packard și ulterior îmbunătățit de Google, Tesseract oferă suport pentru peste 100 de limbi, inclusiv româna și engleza necesare pentru procesarea etichetelor alimentare. Configurarea optimizată pentru texte mici și layout-uri complexe, specifice etichetelor nutriționale, asigură o acuratețe ridicată în extragerea informațiilor.

\subsection{Google Gemini AI}

Modelul de limbaj Gemini a fost ales pentru procesarea și interpretarea listelor de ingrediente datorită performanței superioare în domenii specializate precum chimia și biologia \cite{GptvGemini2024}. Capacitatea multimodală de a procesa simultan text și imagini, alături de accesibilitatea tier-ului gratuit pentru dezvoltare, au făcut din Gemini soluția optimă pentru transformarea ingredientelor cu denumiri științifice în explicații accesibile utilizatorilor.

\subsection{MinIO pentru stocarea obiectelor}

MinIO oferă o soluție de stocare a obiectelor compatibilă cu Amazon S3, optimizată pentru aplicații care gestionează volume mari de imagini. Compatibilitatea S3 permite utilizarea unui ecosistem vast de instrumente și facilitează migrarea viitoare către soluții cloud, în timp ce performanța ridicată și simplitatea deployment-ului ca container Docker se integrează perfect în arhitectura aplicației.

\section{Tehnologii Frontend}

\subsection{React Native}

React Native permite dezvoltarea aplicațiilor mobile native pentru iOS și Android folosind o singură bază de cod JavaScript. Această abordare reduce semnificativ timpul de dezvoltare și costurile de mentenanță, oferind în același timp performanțe apropiate de aplicațiile native prin accesul direct la API-urile platformei.

\subsection{Expo}

Expo simplifică procesul de dezvoltare React Native prin oferirea unui set complet de instrumente și servicii. Gestionarea automată a dependențelor native, sistemul de build în cloud și facilitățile de testare pe device-uri fizice accelerează dezvoltarea și deployment-ul aplicației mobile.

\subsection{React Navigation}

Pentru navigarea între ecrane, React Navigation oferă o soluție robustă și flexibilă, cu suport pentru diverse tipuri de navigare și integrare profundă cu ecosistemul React Native. Gestionarea stării navigației și suportul pentru deep linking facilitează crearea unei experiențe utilizator fluide.

\subsection{Expo Image Picker}

Componenta Expo Image Picker gestionează accesul la camera și galeria foto ale device-ului, oferind o interfață unificată pentru capturarea și selecția imaginilor. Optimizarea automată a imaginilor și suportul pentru diverse formate asigură compatibilitatea între platforme.

\section{Tehnologii de infrastructură}

\subsection{Docker și Docker Compose}

Containerizarea cu Docker asigură portabilitatea aplicației și consistența mediului de rulare între dezvoltare și producție. Docker Compose orchestrează serviciile multiple ale aplicației, simplificând deployment-ul și gestionarea dependențelor între componente.

\subsection{Servicii de autentificare}

Implementarea unui sistem de autentificare bazat pe JWT tokens cu suport pentru refresh tokens asigură securitatea accesului la aplicație. Utilizarea cookie-urilor HTTP-only pentru stocarea token-urilor și implementarea mecanismelor de revoke pentru token-uri oferă un nivel înalt de securitate.

\subsection{Servicii de email}

Integrarea cu fastapi-mail permite trimiterea de email-uri de verificare și resetare parolă, esențiale pentru un sistem de autentificare complet. Configurarea pentru utilizarea serverelor SMTP standard asigură compatibilitatea cu diverse furnizori de servicii email.

\section{Procesarea imaginilor pentru OCR}

Pentru optimizarea rezultatelor OCR, aplicația implementează tehnici de preprocesare a imaginilor folosind Pillow (PIL). Aceste optimizări includ redimensionarea inteligentă pentru păstrarea detaliilor importante, ajustarea contrastului și luminozității pentru îmbunătățirea lizibilității textului și conversia în formate optimizate pentru procesarea ulterioară.

Provocările specifice ale etichetelor alimentare, precum fonturile mici, reflexiile pe suprafețele ambalajelor și layout-urile complexe, necesită algoritmi specializați de preprocesare pentru maximizarea acurateței extragerii textului.

\section{Integrarea AI pentru interpretarea datelor}

Procesarea datelor extrase prin OCR beneficiază de capacitățile avansate ale modelelor de limbaj mari pentru interpretarea și contextualizarea informațiilor nutriționale. Sistemul implementează mecanisme de retry și validare pentru asigurarea consistenței rezultatelor, precum și algoritmi de post-procesare pentru îmbunătățirea calității datelor finale.
